% !TEX root = ../main.tex
\chapter{正文、图表与引用的写法}
\label{chap:guide2}

\section{一段文字只推进一个主要判断}

段落可以按照“主要判断、相关材料、解释、边界”的顺序组织。主要判断告诉读者这一段要说明什么，材料给出依据，解释连接材料与判断，边界则说明哪些情况尚不能推出。并非每一段都必须包含四个句子，但这四个问题可以用来发现论证中的缺口。把多个独立判断塞进同一长句，往往使读者难以确定哪些证据支持哪个结论。

修改文字时，应优先修正含义和证据关系，再处理语气与措辞。将“效果很好”替换为“效果显著”，只是换了一个词，并没有补足指标、条件和依据。表~\ref{tab:guide2} 中的推荐写法保留了未知信息，没有把待完成的检查改写为确定结果。它们是原创教学对照句，不是从实际论文摘录的文字。证据与引用清单可以辅助记录每次修改是否改变了原有结论的强度。\cite{guide2}

\begin{table}[htbp]
  \centering
  \small
  \renewcommand{\arraystretch}{1.25}
  \caption{表述修改的教学对照}
  \label{tab:guide2}
  \begin{tabular}{@{}>{\raggedright\arraybackslash}p{.2\linewidth}>{\raggedright\arraybackslash}p{.52\linewidth}>{\raggedright\arraybackslash}p{.2\linewidth}@{}}
    \toprule
    原始表述 & 推荐表述 & 修改目的 \\
    \midrule
    这个方法效果很好。 & 需要先说明评价指标与比较条件，再描述观察到的变化。 & 补足判断依据 \\
    所有情况都适用。 & 本文仅讨论材料中明确列出的条件，其他情形尚未验证。 & 限定适用范围 \\
    参考文献已经整理完了。 & 已完成条目录入，作者、年份与正文引用仍需逐项核对。 & 区分操作与验证 \\
    \bottomrule
  \end{tabular}
\end{table}

\section{公式需要定义、计算和文字解释}

公式进入正文之前，先说明它回答什么问题，并在首次使用时定义符号。以式~\eqref{eq:mean} 的算术平均数为例，设 $x_i$ 表示第 $i$ 个教学数值，$n$ 表示数值个数，$\bar{x}$ 表示算术平均数，则

\begin{equation}
\bar{x}=\frac{1}{n}\sum_{i=1}^{n}x_i.
\label{eq:mean}
\end{equation}

本例的数据为 2、4、6，均为独立构造的教学数值，单位为任意教学单位。于是 $n=3$，数值之和为 12，均值为 4。这个计算仅说明公式、变量和表述如何对应，不代表任何实际实验的测量结果。正文引用公式时应使用自动编号，不要手工写死编号。公式后的文字还应解释结果能够回答的问题，避免读者只看到一个孤立数值。

同一个符号应尽量保持同一含义。如果后文需要改变单位、集合或统计范围，必须在改变的位置说明。不要仅凭均值判断数据的离散程度，也不要从三个教学数值推导方法的一般有效性。这里保留简单算例，是为了让读者能够独立核算，并检查正文、公式和数据文件是否一致。复杂研究仍需根据实际设计选择指标和报告方式。

\section{图表与引用都应能够回到来源}

表格适合并列比较，图形适合呈现关系或流程。选择形式时，先判断读者需要完成什么阅读任务，再决定是否增加图表。标题应说明图表的对象，表头应给出字段含义，涉及数值时还要交代单位和统计范围。正文需要指出读者应观察什么，而不是只写“如下图所示”。修改图表顺序后，应重新检查标题、标签和正文引用是否仍然对应。

引用用于说明某个定义、方法或判断来自哪里，不能替代作者自己的解释。将条目放进文献数据库之后，仍需检查正文是否确实引用了它，引用位置是否符合所支撑的内容。对于缺失的作者、年份或出版信息，应回到材料核对；无法确认的字段保持缺失，不凭印象补齐。本指南引用的两份说明均随项目提供，属于原创教学文档，并不声称是外部发表的研究成果。
